%=======================================================================
%  IE 5311-001 (D01) Principles of Optimization
%  Homework 06 solutions
%
%  Packages: geometry, amsmath, amssymb, amsthm, booktabs, array.
%  No figure and no external file: this .tex compiles on its own.
%=======================================================================
\documentclass[11pt]{article}

\usepackage[margin=1in]{geometry}
\usepackage{amsmath}
\usepackage{amssymb}
\usepackage{amsthm}
\usepackage{booktabs}
\usepackage{array}

\DeclareMathOperator*{\minimize}{min}
\DeclareMathOperator*{\maximize}{max}
\DeclareMathOperator*{\argmin}{arg\,min}
\DeclareMathOperator{\conv}{conv}
\DeclareMathOperator{\interior}{int}
\DeclareMathOperator{\tr}{tr}
\newcommand{\R}{\mathbb{R}}
\newcommand{\T}{^{\!\top}}

\theoremstyle{definition}
\newtheorem{remark}{Remark}

\setlength{\parskip}{0.4em}
\setlength{\parindent}{0pt}

\title{IE 5311-001 (D01) Principles of Optimization\\
       Homework 6}
\author{Lidivine Kengne}
\date{Fall 2026}

\begin{document}
\maketitle

\section*{Problem 0.\quad[1 extra-credit point]}

I have completed the course evaluation before submitting this homework.

\section*{Problem 1.\quad Maximum entropy from moment data \quad[25 pts]}

\subsection*{1.\quad The formulation}

\paragraph{Sets.} $I = \{1,\dots,n\}$, the faces of the die, indexed by
$i$.

\paragraph{Parameters.}
\begin{center}
\begin{tabular}{@{}ll@{}}
\toprule
Symbol & Meaning\\
\midrule
$n$ & number of faces, a positive integer\\
$m_1$ & estimate of $\sum_{i \in I} i\,p_i$, in face-index units\\
$m_2$ & estimate of $\sum_{i \in I} i^2 p_i$, in squared face-index units\\
\bottomrule
\end{tabular}
\end{center}

\paragraph{Decision variables.} $p_i \ge 0$ for $i \in I$, the
probability of face $i$; dimensionless.

\paragraph{Program.} Maximising entropy is minimising its negative, so
\begin{equation}\label{eq:p1primal}
\begin{aligned}
  \minimize_{p}\quad & -H(p) \;=\; \sum_{i \in I} p_i \ln p_i\\
  \text{s.t.}\quad   & \sum_{i \in I} p_i = 1,\\
                     & \sum_{i \in I} i\, p_i = m_1,\\
                     & \sum_{i \in I} i^2 p_i = m_2,\\
                     & p_i \ge 0, \quad i \in I,
\end{aligned}
\end{equation}
with the usual convention $0\ln 0 = 0$. Each term $t \mapsto t \ln t$
is strictly convex on $t > 0$, so the objective is convex and the three
constraints are affine: \eqref{eq:p1primal} is a convex program.

\paragraph{Classification (slides 8--10).} Constrained; continuous
variables, compact ($n$ of them); objective convex nonlinear and
separable, constraints linear; parameters deterministic once $m_1, m_2$
are estimated; single stage; single player.

\begin{remark}
The moment data enter only through $m_1$ and $m_2$. Replacing the die
faces $i$ by arbitrary values $v_i$ changes nothing below except that
$i$ and $i^2$ become $v_i$ and $v_i^2$.
\end{remark}

\subsection*{2.\quad The Lagrangian}

Attach free multipliers $\nu_0, \nu_1, \nu_2$ to the three equalities.
The nonnegativity is kept in the domain rather than dualised, for the
reason given in Part~5: the logarithm already forces $p > 0$. Then
\begin{equation}\label{eq:p1lag}
  L(p,\nu) \;=\; \sum_{i \in I} p_i \ln p_i
    \;+\; \nu_0\Bigl(\sum_{i \in I} p_i - 1\Bigr)
    \;+\; \nu_1\Bigl(\sum_{i \in I} i\,p_i - m_1\Bigr)
    \;+\; \nu_2\Bigl(\sum_{i \in I} i^2 p_i - m_2\Bigr).
\end{equation}

\subsection*{3.\quad The dual, and its simplification}

$L$ is separable in $p$ and strictly convex, so the infimum is found
termwise. Setting $\partial L / \partial p_i = 0$,
\[
  \ln p_i + 1 + \nu_0 + \nu_1 i + \nu_2 i^2 = 0
  \qquad\Longrightarrow\qquad
  p_i(\nu) = \exp\!\bigl(-1 - \nu_0 - \nu_1 i - \nu_2 i^2\bigr),
\]
which is automatically strictly positive, so the dropped constraint
$p \ge 0$ is satisfied without being imposed. Substituting back, the
first three terms collapse because
$\sum_i p_i \ln p_i = \sum_i p_i(-1-\nu_0-\nu_1 i - \nu_2 i^2)$, leaving
\begin{equation}\label{eq:p1dual}
  g(\nu) \;=\; -\sum_{i \in I}
      \exp\!\bigl(-1-\nu_0-\nu_1 i-\nu_2 i^2\bigr)
      \;-\; \nu_0 - \nu_1 m_1 - \nu_2 m_2 ,
  \qquad
  \maximize_{\nu \in \R^3} \; g(\nu),
\end{equation}
an unconstrained concave maximisation.

\paragraph{Eliminating $\nu_0$.} Write
$Z(\nu_1,\nu_2) = \sum_{i \in I} e^{-\nu_1 i - \nu_2 i^2}$, the
partition function. Then
$g = -e^{-1-\nu_0}Z - \nu_0 - \nu_1 m_1 - \nu_2 m_2$, and
$\partial g/\partial \nu_0 = e^{-1-\nu_0}Z - 1 = 0$ gives
$\nu_0 = \ln Z - 1$. Substituting,
\begin{equation}\label{eq:p1dualred}
  \maximize_{\nu_1,\nu_2 \in \R}\quad
  -\ln Z(\nu_1,\nu_2) - \nu_1 m_1 - \nu_2 m_2
  \qquad\Longleftrightarrow\qquad
  \minimize_{\nu_1,\nu_2 \in \R}\quad
  \ln Z(\nu_1,\nu_2) + \nu_1 m_1 + \nu_2 m_2 .
\end{equation}

This is the payoff of dualising. The primal has $n$ variables and three
equalities; the reduced dual is smooth, unconstrained and has
\emph{two} variables whatever $n$ is --- one per moment constraint. The
primal solution is recovered from the dual optimum as the Gibbs
distribution
\begin{equation}\label{eq:p1gibbs}
  p_i^\star \;=\;
  \frac{\exp\!\bigl(-\nu_1^\star i - \nu_2^\star i^2\bigr)}
       {Z(\nu_1^\star,\nu_2^\star)}, \qquad i \in I .
\end{equation}
Maximum entropy subject to moment constraints therefore always produces
an exponential family in the constraint functions; the multipliers are
its natural parameters.

\subsection*{4.\quad Slater's criterion}

The objective is convex and all three constraints are affine, so Slater
reduces to the existence of a point in the \emph{relative interior} of
the domain that satisfies them: some $p$ with $p_i > 0$ for every
$i \in I$ and
$\sum_i p_i = 1$, $\sum_i i p_i = m_1$, $\sum_i i^2 p_i = m_2$.

As $p$ ranges over the simplex, the pair
$\bigl(\sum_i i p_i, \sum_i i^2 p_i\bigr)$ ranges over
$\conv\{(i,i^2) : i \in I\}$, and $p$ may be taken strictly positive
exactly when the pair lies in the interior of that hull. Hence
\begin{equation}\label{eq:p1slater}
  \text{Slater holds}
  \quad\Longleftrightarrow\quad
  (m_1,m_2) \in \interior \conv\bigl\{(i,i^2) : i \in I\bigr\}.
\end{equation}

The hull is explicit. Its upper boundary is the single chord from
$(1,1)$ to $(n,n^2)$, of slope $n+1$; its lower boundary is the
piecewise-linear interpolation through consecutive points, the piece
over $m_1 \in [j,j+1]$ having slope $2j+1$. So, with
$j = \lfloor m_1 \rfloor$, Slater holds if and only if
\begin{equation}\label{eq:p1slater2}
  1 < m_1 < n
  \qquad\text{and}\qquad
  (2j+1)\,m_1 - j(j+1) \;<\; m_2 \;<\; (n+1)\,m_1 - n .
\end{equation}

Three consequences are worth stating. First, $n \ge 3$ is necessary:
for $n = 2$ the two points $(1,1)$ and $(2,4)$ span a segment with
empty interior, and the three equalities already pin $p$ down. Second,
since the points lie on the convex curve $y = x^2$, the lower bound in
\eqref{eq:p1slater2} exceeds $m_1^2$, so $m_2 > m_1^2$ --- a strictly
positive variance --- is necessary but not sufficient. Third, on the
boundary the feasible $p$ is unique and supported on at most two faces,
so some $p_i = 0$: no strictly positive feasible point exists, Slater
fails, and the dual optimum is not attained, the multipliers receding
to infinity as the corresponding $p_i$ approach zero.

\subsection*{5.\quad KKT conditions}

Attach $\mu_i \ge 0$ to $p_i \ge 0$ and keep $\nu_0,\nu_1,\nu_2$ free.
A pair $(p^\star,\nu^\star,\mu^\star)$ satisfies the KKT conditions when
\[
\begin{array}{llll}
  \text{stationarity} &
    \ln p_i^\star + 1 + \nu_0^\star + \nu_1^\star i
      + \nu_2^\star i^2 - \mu_i^\star = 0, & i \in I,\\[2pt]
  \text{primal feasibility} &
    \displaystyle\sum_{i \in I} p^\star_i = 1,\quad
    \sum_{i \in I} i p^\star_i = m_1,\quad
    \sum_{i \in I} i^2 p^\star_i = m_2,\quad
    p^\star_i \ge 0, & i \in I,\\[8pt]
  \text{dual feasibility} & \mu^\star_i \ge 0, & i \in I,\\[2pt]
  \text{complementary slackness} & \mu^\star_i p^\star_i = 0, & i \in I.
\end{array}
\]

The nonnegativity multipliers are redundant here. Stationarity involves
$\ln p_i^\star$, so it can only hold at $p_i^\star > 0$; complementary
slackness then forces $\mu_i^\star = 0$ for every $i$. The system
collapses to
\[
  p_i^\star = \exp\!\bigl(-1-\nu_0^\star-\nu_1^\star i
      - \nu_2^\star i^2\bigr), \quad i \in I,
  \qquad\text{together with the three moment equations,}
\]
which is \eqref{eq:p1gibbs} with $\nu_0^\star = \ln Z - 1$. The entropy
acts as its own barrier: the constraint $p \ge 0$ is never active. Since
the problem is convex, under \eqref{eq:p1slater} these conditions are
both necessary and sufficient for optimality.

\subsection*{Numerical check}

\texttt{p1\_maxent.py} solves \eqref{eq:p1primal} with an interior-point
conic solver and \eqref{eq:p1dualred} by quasi-Newton on two variables,
for $n = 6$.

\begin{center}
\begin{tabular}{@{}lrrrrrr@{}}
\toprule
case & $m_1$ & $m_2$ & $\nu_1^\star$ & $\nu_2^\star$
  & $-H(p^\star)$ & duality gap\\
\midrule
fair    & 3.5 & 15.1667 & $-0.0000000$ & $+0.0000000$
  & $-1.7917594677$ & $1.5\times10^{-9}$\\
skewed  & 4.0 & 18.0    & $-1.2417518$ & $+0.1429377$
  & $-1.6965357595$ & $2.3\times10^{-9}$\\
tight   & 4.5 & 21.0    & $-5.6053109$ & $+0.6192405$
  & $-1.2665057255$ & $1.4\times10^{-10}$\\
\bottomrule
\end{tabular}
\end{center}

The largest discrepancy between the primal solution and the Gibbs form
\eqref{eq:p1gibbs} rebuilt from the dual optimum was
$1.8 \times 10^{-5}$. The fair case returns $\nu^\star = 0$ and the
uniform distribution, as it must, since the uniform law already matches
both moments and maximises entropy unconstrained.

Taking $m_1 = 4$, $m_2 = 22$ puts the pair on the upper chord, where
\eqref{eq:p1slater2} fails. The feasible set collapses to the single
point $p = (0.4,0,0,0,0,0.6)$ --- supported only on faces $1$ and $6$,
the chord's endpoints --- with $-H(p) = -0.6730$, confirming that
boundary moment data leave no strictly positive feasible distribution.

\section*{Problem 2.\quad Bilevel capacity and production \quad[15 pts]}

\subsection*{1.\quad The bilevel program}

\paragraph{Sets.} $J = \{1,2\}$, the two production facilities, indexed
by $j$.

\paragraph{Parameters.}
\begin{center}
\begin{tabular}{@{}llr@{}}
\toprule
Symbol & Meaning & Value\\
\midrule
$d$ & total demand, in units & $4$\\
$\bar{x}$ & upper bound on installed capacity at facility 1, in units & $4$\\
$c^{\mathrm{cap}}$ & capacity investment cost at facility 1, per unit & $0.8$\\
$c^{\mathrm{prod}}$ & production cost at facility 2, per unit & $2$\\
$q_1, q_2$ & follower cost coefficients, per squared unit
  & $\tfrac12,\; 1$\\
\bottomrule
\end{tabular}
\end{center}

\paragraph{Decision variables.} The leader chooses $x \in [0,\bar{x}]$,
the capacity installed at facility 1. The follower, having observed
$x$, chooses $y_j \ge 0$, the quantity produced at facility $j$.

\paragraph{Program.} The leader minimises its own cost anticipating
that the follower will respond optimally:
\begin{equation}\label{eq:p2bilevel}
\begin{aligned}
  \minimize_{x,\,y}\quad & c^{\mathrm{cap}} x + c^{\mathrm{prod}} y_2\\
  \text{s.t.}\quad & 0 \le x \le \bar{x},\\
  & y \in \argmin_{y'}\Bigl\{\,
      q_1 y_1'^{\,2} + q_2 y_2'^{\,2}
      \;:\; y_1' + y_2' = d,\;\; 0 \le y_1' \le x,\;\; y_2' \ge 0
      \Bigr\}.
\end{aligned}
\end{equation}

With the data, the leader minimises $0.8x + 2y_2$ and the follower
minimises $\tfrac12 y_1^2 + y_2^2$. The leader pays for capacity at
facility 1 and for output at facility 2; the follower pays a convex
production cost and ignores the leader's costs entirely. The two
objectives conflict: the leader wants $y_1$ large to avoid the cost
$c^{\mathrm{prod}} y_2$, but only the capacity $x$ --- which it must
pay for --- forces the follower's hand.

\paragraph{Classification (slides 8--10).} Constrained; continuous
variables, compact; objectives and constraints convex quadratic and
linear; deterministic; two stages in the Stackelberg sense; \emph{two
players}, which is what makes the problem bilevel rather than a single
nonlinear program.

\subsection*{2.\quad Follower KKT, and the single-level reformulation}

For fixed $x$ the follower solves a convex quadratic program. Attach
$\lambda$ (free) to the demand equality, $\mu \ge 0$ to $y_1 \le x$,
and $\alpha_j \ge 0$ to $y_j \ge 0$:
\[
  \mathcal{L}(y,\lambda,\mu,\alpha)
  = \tfrac12 y_1^2 + y_2^2 + \lambda\,(y_1 + y_2 - d)
    + \mu\,(y_1 - x) - \alpha_1 y_1 - \alpha_2 y_2 .
\]
The KKT conditions are
\begin{equation}\label{eq:p2kkt}
\begin{array}{ll}
  \text{stationarity} & y_1 + \lambda + \mu - \alpha_1 = 0, \qquad
                        2y_2 + \lambda - \alpha_2 = 0,\\[2pt]
  \text{primal feasibility} & y_1 + y_2 = d, \quad y_1 \le x, \quad
                              y_1, y_2 \ge 0,\\[2pt]
  \text{dual feasibility} & \mu \ge 0, \quad \alpha_1, \alpha_2 \ge 0,\\[2pt]
  \text{complementary slackness} & \mu\,(y_1 - x) = 0, \quad
      \alpha_1 y_1 = 0, \quad \alpha_2 y_2 = 0 .
\end{array}
\end{equation}

\paragraph{Why they are necessary and sufficient here.} Necessity does
not need Slater: every follower constraint is affine, so the linearity
constraint qualification holds at every feasible point and KKT is
necessary at any minimiser. Sufficiency is the convexity half: the
follower's objective has Hessian $\operatorname{diag}(1,2) \succ 0$ and
the feasible set is a polyhedron, so any KKT point is a global
minimiser. Strict convexity gives more --- the minimiser is
\emph{unique} for each $x$ --- which matters for
\eqref{eq:p2bilevel} itself: with a unique follower response there is
no gap between the optimistic and pessimistic readings of the bilevel
problem, and $\argmin$ is a genuine function of $x$.

Replacing the inner $\argmin$ by \eqref{eq:p2kkt} gives the
single-level problem
\begin{equation}\label{eq:p2mpcc}
\begin{aligned}
  \minimize_{x,y,\lambda,\mu,\alpha}\quad
    & c^{\mathrm{cap}} x + c^{\mathrm{prod}} y_2\\
  \text{s.t.}\quad & 0 \le x \le \bar{x}, \quad
    \text{all conditions in \eqref{eq:p2kkt}} ,
\end{aligned}
\end{equation}
a mathematical program with complementarity constraints. It is not a
convex program: the three complementarity conditions are the only
nonconvexity, and they are what Part~3 removes.

\subsection*{3.\quad Disjunctive linearisation}

Each complementarity pair has the form $ab = 0$ with $a,b \ge 0$, which
is the disjunction $(a = 0) \vee (b = 0)$. Introducing a binary
$z \in \{0,1\}$ and valid upper bounds $M_a, M_b$,
\[
  ab = 0,\; a,b \ge 0
  \qquad\Longleftrightarrow\qquad
  a \le M_a z, \quad b \le M_b (1-z), \quad z \in \{0,1\},
\]
provided $a \le M_a$ and $b \le M_b$ hold at every point of interest.
Applying this to the three pairs of \eqref{eq:p2kkt} gives the mixed
integer quadratically constrained program
\begin{equation}\label{eq:p2mip}
\begin{aligned}
  \minimize\quad & c^{\mathrm{cap}} x + c^{\mathrm{prod}} y_2\\
  \text{s.t.}\quad
    & y_1 + \lambda + \mu - \alpha_1 = 0, \qquad
      2y_2 + \lambda - \alpha_2 = 0,\\
    & y_1 + y_2 = d, \quad y_1 \le x, \quad 0 \le x \le \bar{x},
      \quad y \ge 0, \quad \mu, \alpha \ge 0,\\
    & \mu \le M^{\mathrm{dual}} z_\mu, \qquad
      x - y_1 \le M^{\mathrm{prim}}(1 - z_\mu),\\
    & \alpha_1 \le M^{\mathrm{dual}} z_1, \qquad
      y_1 \le M^{\mathrm{prim}}(1 - z_1),\\
    & \alpha_2 \le M^{\mathrm{dual}} z_2, \qquad
      y_2 \le M^{\mathrm{prim}}(1 - z_2),\\
    & z_\mu, z_1, z_2 \in \{0,1\} .
\end{aligned}
\end{equation}
The objective and all remaining constraints are linear, so
\eqref{eq:p2mip} is a mixed integer \emph{linear} program.

\paragraph{Justifying the bounds.} A big-M is valid only if it excludes
no KKT point; one chosen too small silently replaces the follower's
problem by a different one, and the leader then optimises against a
response that is not the follower's.

\emph{Primal side.} From $y_1 + y_2 = d$ and $y \ge 0$ we get
$y_1, y_2 \in [0,d]$. From $0 \le y_1 \le x \le \bar{x}$ the slack
satisfies $x - y_1 \in [0,\bar{x}]$. With $d = \bar{x} = 4$,
$M^{\mathrm{prim}} = 4$ is valid and attained.

\emph{Dual side.} First, $\alpha_2 = 0$ always. If $\alpha_2 > 0$ then
$y_2 = 0$, hence $y_1 = d = 4$; stationarity then gives
$\lambda = \alpha_2 > 0$ from the second equation and
$4 + \lambda + \mu = 0$ from the first, so $\lambda \le -4$ --- a
contradiction. Therefore $y_2 > 0$ at every follower optimum and
$\lambda = -2y_2 \in [-2d, 0]$. When $y_1 > 0$ we have $\alpha_1 = 0$
and
\[
  \mu = -\lambda - y_1 = 2y_2 - y_1 \;\le\; 2d = 8 ,
\]
so $M^{\mathrm{dual}} = 8$ is valid. It is also tight: eliminating
$y_2$ gives the follower objective
$\tfrac32 y_1^2 - 8y_1 + 16$, minimised at $y_1 = 8/3$, so
$y_1 = \min\{x, 8/3\}$ and, for $x < 8/3$,
$\mu(x) = 8 - 3x$, which reaches $8$ at $x = 0$.

\paragraph{The cost of getting it wrong.} \texttt{p2\_bilevel.py}
re-solves \eqref{eq:p2mip} with $M^{\mathrm{dual}} = 2$. Every
$x < 2$ becomes infeasible, so the model asserts the follower has no
optimal response on more than half the leader's range. The optimum
nevertheless survives, because it sits at $x^\star = 8/3$ where
$\mu^\star = 0$, and the run still returns $4.8$. That is the point
worth taking: an unjustified bound can return the right number while
describing the wrong problem, so the bound has to be argued, not tuned
until the answer looks plausible.

\paragraph{Alternatives.} The disjunction can also be imposed directly,
without any $M$: as an SOS1 constraint on each pair, or as indicator
constraints $z = 1 \Rightarrow a = 0$. Both avoid the validity question
and the weak continuous relaxation that a large $M$ produces, at the
cost of a branching rule the solver must support.

\subsection*{Solution and numerical check}

Since $y_1 = \min\{x, 8/3\}$ and $y_2 = d - y_1$, the leader's cost is
piecewise linear in $x$:
\[
  c^{\mathrm{cap}} x + c^{\mathrm{prod}} y_2(x) =
  \begin{cases}
    8 - 1.2\,x, & 0 \le x < 8/3 \quad \text{(decreasing)},\\[2pt]
    0.8\,x + 8/3, & 8/3 \le x \le 4 \quad \text{(increasing)},
  \end{cases}
\]
so the optimum is the kink:
\[
  x^\star = \tfrac83 \approx 2.6667, \qquad
  y^\star = \bigl(\tfrac83, \tfrac43\bigr), \qquad
  \lambda^\star = -\tfrac83, \qquad
  \mu^\star = \alpha_1^\star = \alpha_2^\star = 0,
\]
with leader cost $0.8 \cdot \tfrac83 + 2 \cdot \tfrac43 = 4.8$. A grid
search over $400{,}001$ values of $x$ with the follower solved exactly
returns $x^\star = 2.666670$ and cost $4.800003$; the MILP
\eqref{eq:p2mip} returns $x^\star = 2.66666667$ and cost
$4.80000000$.

\begin{remark}
The optimum is degenerate: $y_1^\star = x^\star$ makes the capacity
constraint active while $\mu^\star = 0$, so \emph{both} sides of that
complementarity vanish and both branches $z_\mu = 0$ and $z_\mu = 1$
are feasible there. This is the weak-complementarity case, and it is
also why the invalid bound above escaped detection.
\end{remark}

\section*{Problem 3.\quad Two-stage stochastic LP by Benders \quad[20 pts]}

\subsection*{1.\quad The two-stage stochastic linear program}

\paragraph{Sets.} $I = \{1,2,3\}$, the assets, indexed by $i$;
$S = \{1,2,3\}$, the scenarios, indexed by $s$.

\paragraph{Parameters.}
\begin{center}
\begin{tabular}{@{}lll@{}}
\toprule
Symbol & Meaning & Value\\
\midrule
$c_i$ & transaction cost per unit of wealth placed in asset $i$
  & $(0.004,\,0.006,\,0.003)$\\
$p_s$ & probability of scenario $s$ & $(0.3,\,0.4,\,0.3)$\\
$G_{si}$ & gross return $1+r_{si}$ of asset $i$ in scenario $s$
  & matrix below\\
$L_s$ & required liability in scenario $s$, in units of wealth
  & $(1.01,\,1.00,\,1.02)$\\
$k_s$ & borrowing cost coefficient in scenario $s$, per unit borrowed
  & $(0.08,\,0.09,\,0.10)$\\
\bottomrule
\end{tabular}
\end{center}
\[
  G = \begin{bmatrix}
    1.06 & 1.01 & 0.98\\
    0.95 & 1.03 & 1.02\\
    1.10 & 0.97 & 0.96
  \end{bmatrix}.
\]

\paragraph{Decision variables.} First stage, before the scenario is
known: $x_i \ge 0$, the fraction of initial wealth placed in asset $i$.
Second stage, after scenario $s$ is observed: $b_s \ge 0$, the amount
borrowed in scenario $s$, in units of wealth.

\paragraph{Program.} In extensive form,
\begin{equation}\label{eq:p3ef}
\begin{aligned}
  \minimize_{x,\,b}\quad
    & \sum_{i \in I} c_i x_i
      + \sum_{s \in S} p_s k_s b_s\\
  \text{s.t.}\quad
    & \sum_{i \in I} x_i = 1,\\
    & \sum_{i \in I} G_{si} x_i + b_s \;\ge\; L_s, & s \in S,\\
    & x_i \ge 0,\; i \in I, \qquad b_s \ge 0,\; s \in S .
\end{aligned}
\end{equation}
Equivalently, in the two-stage form that Benders exploits,
\begin{equation}\label{eq:p3two}
  \minimize_{x}\;\Bigl\{\,
    \sum_{i \in I} c_i x_i + \sum_{s \in S} p_s\, Q_s(x)
    \;:\; \sum_{i \in I} x_i = 1,\;\; x \ge 0 \Bigr\},
  \qquad
  Q_s(x) = \minimize_{b_s \ge 0}\;\Bigl\{\,
    k_s b_s : b_s \ge L_s - \textstyle\sum_{i} G_{si} x_i \Bigr\}.
\end{equation}

\paragraph{Classification (slides 8--10).} Constrained; continuous
variables, compact; linear objective and constraints; \emph{stochastic}
parameters on a finite support of three scenarios; two decision stages;
single player.

\subsection*{2.\quad The second-stage subproblem and its dual}

Fix $x$ and write the shortfall $h_s(x) = L_s - \sum_i G_{si} x_i$. The
scenario subproblem and its dual, with $\pi_s$ attached to the
liability row, are
\[
  Q_s(x) = \minimize_{b_s \ge 0}\; k_s b_s
    \;\;\text{s.t.}\;\; b_s \ge h_s(x) \;\;[\pi_s \ge 0]
  \qquad\Longleftrightarrow\qquad
  Q_s(x) = \maximize_{\pi_s}\; \pi_s h_s(x)
    \;\;\text{s.t.}\;\; 0 \le \pi_s \le k_s .
\]
The upper bound comes from the sign condition on $b_s$: the Lagrangian
is $(k_s - \pi_s) b_s + \pi_s h_s(x)$, whose infimum over $b_s \ge 0$
is $-\infty$ unless $\pi_s \le k_s$. Both are solved in closed form,
\begin{equation}\label{eq:p3sub}
  Q_s(x) = k_s \bigl[h_s(x)\bigr]_+,
  \qquad
  \pi_s^\star \in
  \begin{cases}
    \{k_s\}, & h_s(x) > 0,\\
    [\,0,\,k_s\,], & h_s(x) = 0,\\
    \{0\}, & h_s(x) < 0 .
  \end{cases}
\end{equation}

Two structural facts follow, and both are used in Part~3. The
subproblem is feasible for \emph{every} $x$, since $b_s$ may be taken
arbitrarily large --- the problem has relatively complete recourse, so
no feasibility cuts are ever needed. And the dual feasible set
$[0,k_s]$ does not depend on $x$: this is fixed recourse, which is what
makes a cut generated at one $x$ valid at all others.

\subsection*{3.\quad The Benders master problem and its optimality cuts}

Let $\theta$ represent $\sum_s p_s Q_s(x)$. For any dual-feasible
$\pi^{(t)}$, weak duality in \eqref{eq:p3sub} gives
$Q_s(x) \ge \pi_s^{(t)} h_s(x)$ \emph{for every} $x$, since
$\pi^{(t)}$ stays feasible as $x$ moves. Taking the expectation,
\begin{equation}\label{eq:p3cut}
  \theta \;\ge\; \sum_{s \in S} p_s \pi_s^{(t)}
    \Bigl( L_s - \sum_{i \in I} G_{si} x_i \Bigr)
  \;=\; \alpha_t - \sum_{i \in I} \beta_{ti} x_i,
  \qquad
  \alpha_t = \sum_{s} p_s \pi^{(t)}_s L_s, \quad
  \beta_{ti} = \sum_{s} p_s \pi^{(t)}_s G_{si},
\end{equation}
a globally valid optimality cut. After $T$ iterations the master is
\begin{equation}\label{eq:p3master}
\begin{aligned}
  \minimize_{x,\,\theta}\quad & \sum_{i \in I} c_i x_i + \theta\\
  \text{s.t.}\quad & \sum_{i \in I} x_i = 1, \quad x \ge 0,
    \quad \theta \ge 0,\\
  & \theta \ge \alpha_t - \sum_{i \in I} \beta_{ti} x_i,
    & t = 1,\dots,T .
\end{aligned}
\end{equation}
The bound $\theta \ge 0$ keeps the first master bounded and is valid
because $Q_s \ge 0$. Solving \eqref{eq:p3master} gives a lower bound
$\mathrm{LB} = c\T x^{(t)} + \theta^{(t)}$, since the master is a
relaxation; evaluating \eqref{eq:p3sub} at $x^{(t)}$ gives a feasible
point and hence
$\mathrm{UB} = \min_{\tau \le t}\{c\T x^{(\tau)} + \sum_s p_s
Q_s(x^{(\tau)})\}$. \textbf{Convergence criterion:} stop when
$\mathrm{UB} - \mathrm{LB} \le \varepsilon$. The method terminates
finitely because only finitely many distinct $\pi$ vectors can be
generated --- the dual feasible region is a box with $2^{|S|}$ vertices
--- so only finitely many distinct cuts exist.

\subsection*{4.\quad Implementation and results}

\texttt{p3\_benders.py} implements \eqref{eq:p3master} with
$\varepsilon = 10^{-9}$ and solves \eqref{eq:p3ef} separately as a
reference.

\begin{center}
\begin{tabular}{@{}rrrr@{}}
\toprule
iteration & lower bound & upper bound & gap\\
\midrule
1 & 0.0030000000 & 0.0055200000 & $2.52 \times 10^{-3}$\\
2 & 0.0034117647 & 0.0038000000 & $3.88 \times 10^{-4}$\\
3 & 0.0036428571 & 0.0038000000 & $1.57 \times 10^{-4}$\\
4 & 0.0037885714 & 0.0037885714 & $3.90 \times 10^{-18}$\\
\bottomrule
\end{tabular}
\end{center}

\paragraph{Optimal solution.}
\[
  x^\star = \bigl(\tfrac37,\, 0,\, \tfrac47\bigr)
          = (0.428571,\, 0,\, 0.571429), \qquad
  b^\star = (0,\, 0.01,\, 0), \qquad
  z^\star = 0.0037885714 ,
\]
reached in \textbf{4 iterations}. The extensive form \eqref{eq:p3ef}
returns the same objective to $2.2 \times 10^{-18}$ and the same
portfolio to $1.2 \times 10^{-15}$.

\begin{center}
\begin{tabular}{@{}rrrrrr@{}}
\toprule
scenario & wealth $\sum_i G_{si}x^\star_i$ & liability $L_s$
  & shortfall & $b_s^\star$ & $\pi_s^\star$\\
\midrule
1 & 1.014286 & 1.01 & $-4.29\times10^{-3}$ & 0 & 0\\
2 & 0.990000 & 1.00 & $+1.00\times10^{-2}$ & 0.01 & 0.09\\
3 & 1.020000 & 1.02 & $0$ & 0 & $[0,\,0.10]$\\
\bottomrule
\end{tabular}
\end{center}

\paragraph{Reading the solution.} Asset 2 is never held: it carries the
highest transaction cost, $0.006$, and is never the best return in any
scenario, so it is dominated. The remaining two are mixed so that
scenario 3 --- the one with the largest liability, $1.02$ --- is met
\emph{exactly}, scenario 1 is met with slack, and only scenario 2
requires borrowing, of $0.01$ at coefficient $0.09$ with probability
$0.4$. The expected borrowing cost, $0.00036$, is about one tenth of
the total; the rest is the transaction cost $0.024/7 = 0.003429$.

\begin{remark}
Scenario 3 is tight with $b_3^\star = 0$, so by \eqref{eq:p3sub} its
subproblem dual is \emph{not unique}: any $\pi_3 \in [0,0.10]$ is
optimal, and the extensive form reports $0.083810$, strictly interior.
This threatens nothing. The cut \eqref{eq:p3cut} is valid for any
dual-feasible $\pi$, not only an optimal one, so a degenerate tie can
weaken a cut but cannot invalidate it. Re-running the method with
$\pi_3/k_3 \in \{0,\,0.25,\,0.5,\,0.838,\,1\}$ gave 4 iterations and
the same optimum in every case.
\end{remark}

\section*{Problem 4.\quad Tracking error: QP, SOCP and SDP \quad[30 pts]}

\subsection*{1.\quad The convex quadratic program}

\paragraph{Sets.} $I = \{1,2,3\}$, the assets, indexed by $i$.

\paragraph{Parameters.} $\bar{x} = (0.7,\,0.2,\,0.1)$, the benchmark
weights, dimensionless; $\Sigma = \operatorname{diag}(4,1,1)$, the
return covariance matrix; $\rho = 1$, the variance budget.

\paragraph{Decision variables.} $x_i \ge 0$, the portfolio weight on
asset $i$, dimensionless.

\paragraph{Program.}
\begin{equation}\label{eq:p4qp}
\begin{aligned}
  \minimize_{x}\quad & \sum_{i \in I}(x_i - \bar{x}_i)^2
    \;=\; \|x - \bar{x}\|_2^2\\
  \text{s.t.}\quad & x\T \Sigma x \le \rho,\\
    & \sum_{i \in I} x_i = 1,\\
    & x_i \ge 0, \quad i \in I .
\end{aligned}
\end{equation}

\paragraph{Why it is convex.} The objective
$\|x-\bar{x}\|_2^2$ has Hessian $2I \succ 0$, so it is strictly convex.
The risk function $x \mapsto x\T\Sigma x$ has Hessian
$2\Sigma = \operatorname{diag}(8,2,2) \succeq 0$, so it is a convex
quadratic and its sublevel set $\{x : x\T\Sigma x \le \rho\}$ is a
convex set --- a solid ellipsoid. The budget equality and the sign
conditions are affine. A convex objective over the intersection of
convex sets is a convex program, and here strict convexity of the
objective makes the minimiser unique.

\paragraph{Classification (slides 8--10).} Constrained; continuous
variables, compact; objective convex quadratic, constraints one convex
quadratic and the rest linear; deterministic; single stage; single
player.

\begin{remark}
The benchmark is itself infeasible:
$\bar{x}\T\Sigma\bar{x} = 4(0.49)+0.04+0.01 = 2.01 > \rho$. The risk
constraint therefore binds at the optimum, and the problem is a genuine
trade-off rather than a trivial one returning $x^\star = \bar{x}$.
Asset 1 carries four times the variance of the others, so matching the
benchmark's $0.7$ weight on it is what the budget forbids.
\end{remark}

\subsection*{2.\quad Second-order cone form with a linear objective}

Write $\mathcal{Q}^{m+1} = \{(u,v) \in \R \times \R^{m} :
\|v\|_2 \le u\}$ for the second-order cone, and
$\Sigma^{1/2} = \operatorname{diag}(2,1,1)$, so that
$x\T\Sigma x = \|\Sigma^{1/2}x\|_2^2$. Minimising
$\|x-\bar{x}\|_2$ and minimising its square have the same argument, so
introducing an epigraph variable $t$,
\begin{equation}\label{eq:p4socp}
\begin{aligned}
  \minimize_{t,\,x}\quad & t\\
  \text{s.t.}\quad
    & (t,\; x - \bar{x}) \in \mathcal{Q}^{4},\\
    & (1,\; \Sigma^{1/2} x) \in \mathcal{Q}^{4},\\
    & \textstyle\sum_{i \in I} x_i = 1, \qquad x \ge 0 .
\end{aligned}
\end{equation}
The objective is now linear and the only nonlinearity sits in two
second-order cones. The optimal values are related by
$t^\star = \sqrt{\text{optimal value of } \eqref{eq:p4qp}}$, and the
optimal $x$ is the same.

\subsection*{3.\quad The conic dual}

Assign $(\sigma,u) \in \mathcal{Q}^4$ to the first cone,
$(\tau,w) \in \mathcal{Q}^4$ to the second --- the second-order cone is
self-dual, so the dual variables live in the same cone --- together
with $\nu \in \R$ for the budget equality and $z \in \R^3_+$ for
$x \ge 0$. The Lagrangian is
\[
  \begin{aligned}
    L &= t - \bigl[\sigma t + u\T(x-\bar{x})\bigr]
          - \bigl[\tau + w\T \Sigma^{1/2} x\bigr]
          - \nu\Bigl(\sum_i x_i - 1\Bigr) - z\T x\\
      &= t(1-\sigma)
          - x\T\bigl(u + \Sigma^{1/2}w + \nu \mathbf{1} + z\bigr)
          + u\T\bar{x} - \tau + \nu .
  \end{aligned}
\]
The infimum over $t$ is finite only if $\sigma = 1$, and over $x$ only
if $u + \Sigma^{1/2}w + \nu\mathbf{1} + z = 0$. Eliminating
$z = -(u + \Sigma^{1/2}w + \nu\mathbf{1}) \ge 0$ turns that into an
inequality, and $\sigma = 1$ with $(\sigma,u) \in \mathcal{Q}^4$ reads
$\|u\|_2 \le 1$. The dual is
\begin{equation}\label{eq:p4socpdual}
\begin{aligned}
  \maximize_{u,\,w,\,\nu,\,\tau}\quad & \bar{x}\T u + \nu - \tau\\
  \text{s.t.}\quad & \|u\|_2 \le 1, \qquad \|w\|_2 \le \tau,\\
    & u + \Sigma^{1/2} w + \nu \mathbf{1} \le 0 .
\end{aligned}
\end{equation}
Since $\tau$ enters the objective only as $-\tau$ and is otherwise
constrained only by $\|w\|_2 \le \tau$, it is optimal to set
$\tau = \|w\|_2$, giving the simplified dual
\begin{equation}\label{eq:p4socpdual2}
  \maximize_{u,\,w,\,\nu}\quad
    \bar{x}\T u + \nu - \|w\|_2
  \qquad\text{s.t.}\qquad
    \|u\|_2 \le 1, \qquad
    u + \Sigma^{1/2} w + \nu \mathbf{1} \le 0 .
\end{equation}

\paragraph{Slater and strong duality.} For a conic program, Slater asks
for a feasible point strictly inside the cones, the affine equality
being required only in the relative interior sense. Take
$x = (0.1,\,0.45,\,0.45)$ and any $t > \|x-\bar{x}\|_2$. Then
$\sum_i x_i = 1$, $x > 0$ componentwise, and
\[
  \|\Sigma^{1/2}x\|_2^2 = 4(0.1)^2 + (0.45)^2 + (0.45)^2
  = 0.04 + 0.405 = 0.445 < 1 ,
\]
so both cone constraints hold strictly. Slater's condition is therefore
satisfied, and strong duality holds with the dual optimum
\emph{attained}: the optimal values of \eqref{eq:p4socp} and
\eqref{eq:p4socpdual2} coincide and there is no duality gap.

\subsection*{4.\quad Schur complement and the SDP form}

The Schur complement theorem states that for $A \succ 0$,
\[
  \begin{bmatrix} c & b\T \\ b & A \end{bmatrix} \succeq 0
  \qquad\Longleftrightarrow\qquad
  c - b\T A^{-1} b \ge 0 .
\]
Applying it with $A = I_3$ to each quadratic constraint turns both into
linear matrix inequalities. For the objective, introducing
$\tau \ge \|x - \bar{x}\|_2^2$,
\[
  \begin{bmatrix} \tau & (x-\bar{x})\T \\ x-\bar{x} & I_3 \end{bmatrix}
  \succeq 0
  \quad\Longleftrightarrow\quad
  \tau - (x-\bar{x})\T(x-\bar{x}) \ge 0 ,
\]
and for the risk constraint, writing
$x\T\Sigma x = (\Sigma^{1/2}x)\T(\Sigma^{1/2}x)$,
\[
  \begin{bmatrix} 1 & (\Sigma^{1/2}x)\T \\ \Sigma^{1/2}x & I_3
  \end{bmatrix} \succeq 0
  \quad\Longleftrightarrow\quad
  1 - x\T \Sigma x \ge 0 .
\]
Hence the semidefinite program
\begin{equation}\label{eq:p4sdp}
\begin{aligned}
  \minimize_{\tau,\,x}\quad & \tau\\
  \text{s.t.}\quad
    & M_1(\tau,x) :=
      \begin{bmatrix} \tau & (x-\bar{x})\T \\ x-\bar{x} & I_3
      \end{bmatrix} \succeq 0,\\
    & M_2(x) :=
      \begin{bmatrix} 1 & (\Sigma^{1/2}x)\T \\ \Sigma^{1/2}x & I_3
      \end{bmatrix} \succeq 0,\\
    & \textstyle\sum_{i \in I} x_i = 1, \qquad x \ge 0 ,
\end{aligned}
\end{equation}
with a linear objective and affine matrix constraints. Note that
$\tau$ is the \emph{squared} tracking error, so
$\tau^\star = (t^\star)^2$ with $t^\star$ from \eqref{eq:p4socp}; the
two formulations share the same $x^\star$ but report objectives on
different scales, which Part~6 checks explicitly.

\subsection*{5.\quad The dual SDP}

Attach $Y_1 \succeq 0$ and $Y_2 \succeq 0$ to the two blocks, $\nu$ to
the budget and $z \ge 0$ to $x \ge 0$, and partition conformably,
\[
  Y_1 = \begin{bmatrix} \eta_1 & y_1\T \\ y_1 & \Upsilon_1\end{bmatrix},
  \qquad
  Y_2 = \begin{bmatrix} \eta_2 & y_2\T \\ y_2 & \Upsilon_2\end{bmatrix},
  \qquad \eta_j \in \R,\; y_j \in \R^3 .
\]
Using $\langle Y_1, M_1\rangle = \eta_1 \tau
+ 2y_1\T(x-\bar{x}) + \tr \Upsilon_1$ and
$\langle Y_2, M_2\rangle = \eta_2 + 2y_2\T\Sigma^{1/2}x + \tr\Upsilon_2$,
the Lagrangian is
\[
  L = \tau(1 - \eta_1)
    - x\T\bigl(2y_1 + 2\Sigma^{1/2}y_2 + \nu\mathbf{1} + z\bigr)
    + 2y_1\T\bar{x} - \tr\Upsilon_1 - \eta_2 - \tr\Upsilon_2 + \nu ,
\]
so the infimum is finite only when $\eta_1 = 1$ and
$2y_1 + 2\Sigma^{1/2}y_2 + \nu\mathbf{1} + z = 0$, giving
\begin{equation}\label{eq:p4sdpdual}
\begin{aligned}
  \maximize\quad & 2\bar{x}\T y_1 - \tr\Upsilon_1
    - \eta_2 - \tr\Upsilon_2 + \nu\\
  \text{s.t.}\quad
    & \begin{bmatrix} 1 & y_1\T\\ y_1 & \Upsilon_1\end{bmatrix}
      \succeq 0, \qquad
      \begin{bmatrix} \eta_2 & y_2\T\\ y_2 & \Upsilon_2\end{bmatrix}
      \succeq 0,\\
    & 2y_1 + 2\Sigma^{1/2}y_2 + \nu\mathbf{1} \le 0 .
\end{aligned}
\end{equation}
The matrix variables can be eliminated. By Schur, $Y_1 \succeq 0$ with
$\eta_1 = 1$ means $\Upsilon_1 \succeq y_1 y_1\T$, and since
$-\tr\Upsilon_1$ is maximised by taking $\Upsilon_1$ as small as
possible, $\Upsilon_1 = y_1 y_1\T$ and $\tr\Upsilon_1 = \|y_1\|_2^2$.
Likewise $\Upsilon_2 = y_2 y_2\T/\eta_2$ gives
$-\eta_2 - \|y_2\|_2^2/\eta_2$, maximised over $\eta_2 > 0$ at
$\eta_2 = \|y_2\|_2$, worth $-2\|y_2\|_2$. Substituting $a = 2y_1$ and
$w = 2y_2$,
\begin{equation}\label{eq:p4sdpdual2}
  \maximize_{a,\,w,\,\nu}\quad
    \bar{x}\T a - \tfrac14\|a\|_2^2 - \|w\|_2 + \nu
  \qquad\text{s.t.}\qquad
    a + \Sigma^{1/2} w + \nu \mathbf{1} \le 0 .
\end{equation}

\begin{remark}
Comparing \eqref{eq:p4socpdual2} with \eqref{eq:p4sdpdual2}, the two
duals share the constraint and the terms $\nu$ and $-\|w\|_2$, and
differ only in how the tracking error is priced: the SOCP dual carries
the constraint $\|u\|_2 \le 1$, the SDP dual the quadratic
$-\tfrac14\|a\|_2^2$. That is exactly the difference between the
Fenchel conjugates of $\|\cdot\|_2$ and $\|\cdot\|_2^2$, which is the
only place the two primals differ.
\end{remark}

\subsection*{6.\quad Implementation}

\texttt{p4\_conic.py} builds \eqref{eq:p4socp} and \eqref{eq:p4sdp} in
CVXPY and solves each with the open-source conic solvers Clarabel and
SCS; times are the median of nine solves.

\begin{center}
\begin{tabular}{@{}llrrrrr@{}}
\toprule
form & solver & objective & $x_1$ & $x_2$ & $x_3$ & time (ms)\\
\midrule
QP   & Clarabel & 0.0855492535 & 0.461424 & 0.310037 & 0.228539 & ---\\
SOCP & Clarabel & 0.2924880336 & 0.461424 & 0.310037 & 0.228539 & 1.06\\
SOCP & SCS      & 0.2924833119 & 0.461426 & 0.310037 & 0.228538 & 1.16\\
SDP  & Clarabel & 0.0855492525 & 0.461424 & 0.310037 & 0.228539 & 2.53\\
SDP  & SCS      & 0.0855451999 & 0.461424 & 0.310038 & 0.228539 & 1.71\\
\bottomrule
\end{tabular}
\end{center}

\paragraph{Consistency.} The SOCP reports the norm and the SDP its
square, so the check is $(t^\star)^2 = \tau^\star$:
\[
  (t^\star)^2 = 0.0855492498, \qquad
  \tau^\star = 0.0855492525, \qquad
  \bigl|(t^\star)^2 - \tau^\star\bigr| = 2.7\times10^{-9},
\]
against the QP value $0.0855492535$, and the portfolios agree to
$\max_i|x^{\mathrm{socp}}_i - x^{\mathrm{sdp}}_i| =
2.0\times10^{-7}$. The budget holds to $10^{-10}$ and the risk
constraint is active at $x^\star{}\T\Sigma x^\star = 1.0000000048$,
confirming the remark in Part~1. The two formulations describe the same
problem.

\paragraph{The hand-derived duals.} Solving the two duals
numerically and comparing each against its own primal,
\[
  \begin{array}{lrrr}
    \toprule
    \text{dual} & \text{dual value} & \text{primal value}
      & \text{gap}\\
    \midrule
    \text{SOCP, \eqref{eq:p4socpdual2}}
      & 0.2924880247 & 0.2924880336 & 8.8\times10^{-9}\\
    \text{SDP, \eqref{eq:p4sdpdual2}}
      & 0.0855492521 & 0.0855492525 & 3.9\times10^{-10}\\
    \bottomrule
  \end{array}
\]
The vanishing gaps confirm both the derivations and the Slater
argument of Part~3.

\paragraph{Computational performance.} At $n = 3$ the times above are
dominated by solver setup, so they are not a basis for comparison. The
difference is structural and appears as the problem grows: the SDP
carries two $(n+1)\times(n+1)$ blocks where the SOCP carries two cones
of length $n+1$. Repeating the experiment on random instances:

\begin{center}
\begin{tabular}{@{}rrrrr@{}}
\toprule
$n$ & SOCP solve (ms) & SDP solve (ms) & ratio
  & $|(t^\star)^2 - \tau^\star|$\\
\midrule
  3 & 0.12 &   0.57 & $4.8\times$   & $1.9\times10^{-10}$\\
 10 & 0.21 &   1.92 & $9.1\times$   & $6.2\times10^{-11}$\\
 30 & 0.40 &  10.03 & $25.1\times$  & $7.0\times10^{-11}$\\
 60 & 0.85 &  66.54 & $78.7\times$  & $1.4\times10^{-10}$\\
100 & 1.36 & 337.52 & $248.7\times$ & $4.3\times10^{-10}$\\
\bottomrule
\end{tabular}
\end{center}

The two agree to better than $10^{-9}$ throughout while the cost ratio
grows from $5$ to $249$. The conclusion is the practical one: the Schur
complement is the right tool for \emph{establishing} that a quadratic
constraint is conic-representable, but lifting a second-order cone into
a semidefinite block is a strictly worse way to solve the resulting
problem. Use the smallest cone that represents the constraint.

\vspace{1em}
\hrule
\vspace{0.6em}
{\small \sloppy \texttt{p1\_maxent.py} solves the entropy primal and its
two-variable dual and checks the Gibbs recovery and the Slater
boundary; \texttt{p2\_bilevel.py} solves the bilevel problem by grid
search and by the big-M MILP and measures the multiplier range that
justifies the bounds; \texttt{p3\_benders.py} runs Benders against the
extensive form and tests the degenerate dual tie-break;
\texttt{p4\_conic.py} builds the SOCP and SDP in CVXPY, checks both
hand-derived duals, and measures how the two scale.\par}

\end{document}
